\chapter{Complexity Classes}

Given what we have already seen, what is the best time complexity for a problem? When can we hope 
to come up with a \emph{fast} algorithm?

This is exactly what complexity classes are about. While the classes themselves are defined rigorously, 
there are some very important things which we have not yet proved -- they remain open problems and we 
believe them to be true in practice$\dots$

\section{Decision Problems}
Complexity classes are traditionally defined using \emph{decision problems}. A decision problem is a 
computational problem whose answer is either \texttt{YES} or \texttt{NO}.

\paragraph{Examples}
\begin{itemize}
\item Given a graph $G$, is $G$ connected?
\item Given an integer $m$, is $m$ prime?
\item Given a graph $G$ and an integer $k$, does $G$ contain a clique\footnote{Refer \autoref{clique}} of 
size at least $k$?
\end{itemize}

Although not all computational tasks are naturally decision problems, most can be 
reformulated as one.

\section{$P$}

The class $P$ consists of all decision problems that can be solved by a 
deterministic algorithm in at most polynomial time. Let Q be a decision problem and T(Q) be the best
time complexity\footnote{This time complexity is in terms of the input length 
and not the input number itself -- a number $n$ when taken as input is represented by 
a bit string of size $m=\Theta(\log n)$. So an algorithm which takes $O(n)$ time will be considered 
as an exponential time algorithm since $\Theta(n) = \Theta(2^m)$. This is why the normal method of 
trying all numbers from $1$ to $\sqrt n$ will not prove that primality testing of a number is in $P$.} 
of a deterministic algorithm for Q. Then $T(Q) \le \Theta(n^c)$ for some constant $c\ge1$.

\[
P \; = \; \bigcup_{c\ge1} \;\; \{Q: T(Q) \le \Theta(n^c)\}
\]

\begin{center}
\begin{tabular}{|c|c|c|}
\hline
\textbf{Decision Problem(Q)} & \textbf{Algorithm} & \textbf{Runtime (T(Q))} \\
\hline
Primality Test & AKS & $\Theta(m^{12})$ \\
Graph Connectivity & BFS & $\Theta(V+E)$ \\
\hline
\end{tabular}
\end{center}

\section{NP}

The class $NP$ consists of all decision problems whose answer when "YES" can be verified in polynomial 
time. The algorithm which solved the problem outputs a \emph{certificate} along with "YES" -- this 
certificate is information which can be used to verify in polynomial time that the answer is indeed "YES". 

\begin{center}
\begin{tabular}{|c|c|c|}
\hline
\textbf{Decision Problem (Q)} & \textbf{Certificate} & \textbf{Verification Time} \\
\hline
Clique of size $k$ & Set of $k$ vertices & $\Theta(k^2)$ \\

Graph Coloring of k colors & Vertex\;:\;Color mapping & $\Theta(V+E)$ \\
\hline
\end{tabular}
\end{center}

An obvious conclusion from the definition of $NP$ is that $P\subseteq NP$ because any problem that can be 
solved efficiently can also be verified efficiently.

\subsection{The Question of \texorpdfstring{$P$}{P} versus \texorpdfstring{$NP$}{NP}}

Every problem that can be solved efficiently can certainly be verified efficiently. Therefore,

\[
P\subseteq NP.
\]

The central open problem in theoretical computer science is whether this inclusion is strict.

\paragraph{The $P$ versus $NP$ Problem.}

Is

\[
P=NP;?
\]

No proof is currently known.

If $P=NP$, then every problem whose solutions can be verified efficiently can also be solved efficiently.

\section{The Class \texorpdfstring{$co$-$NP$}{co-NP}}

The class $co$-$NP$ contains complements of problems in $NP$.

\paragraph{Definition (Class $co$-$NP$).}
A decision problem belongs to $co$-$NP$ if its complement belongs to $NP$.

Equivalently, a \texttt{NO} answer admits a polynomially verifiable certificate.

For example, the complement of the Clique problem asks whether a graph contains no clique of size at least $k$.

We currently do not know whether

\[
NP=co\text{-}NP.
\]

\section{Polynomial-Time Reductions}

To compare the difficulty of computational problems, we use reductions.

\paragraph{Definition (Polynomial-Time Reduction).}
A problem $A$ is polynomial-time reducible to a problem $B$, written

\[
A\le_p B,
\]

if every instance of $A$ can be transformed into an equivalent instance of $B$ in polynomial time.

Intuitively, if $A\le_p B$ and we know how to solve $B$ efficiently, then we can also solve $A$ efficiently.

Reductions allow us to transfer algorithmic techniques as well as hardness results between problems.

\section{NP-Complete Problems}

Some problems are at least as hard as every problem in $NP$.

\paragraph{Definition (NP-Complete).}
A problem $L$ is NP-complete if

\begin{enumerate}
\item $L\in NP$,
\item For every problem $A\in NP$,
\[
A\le_p L.
\]
\end{enumerate}

Thus every problem in $NP$ can be transformed into $L$ using only polynomial time.

\paragraph{Theorem.}

If any NP-complete problem belongs to $P$, then

\[
P=NP.
\]

\paragraph{Examples of NP-Complete Problems.}

\begin{itemize}
\item SAT
\item 3-SAT
\item Clique
\item Vertex Cover
\item Hamiltonian Cycle
\item Subset Sum
\end{itemize}

The discovery of NP-completeness provides strong evidence that these problems are unlikely to possess polynomial-time algorithms.

\section{Randomized Complexity Classes}

We now introduce complexity classes that permit algorithms to use random bits during computation.

\subsection{The Class \texorpdfstring{$RP$}{RP}}

\paragraph{Definition (Randomized Polynomial Time).}

A decision problem belongs to $RP$ if there exists a randomized polynomial-time algorithm satisfying:

\begin{itemize}
\item If the correct answer is \texttt{NO}, the algorithm always outputs \texttt{NO}.
\item If the correct answer is \texttt{YES}, the algorithm outputs \texttt{YES} with probability at least $\frac12$.
\end{itemize}

Thus an $RP$ algorithm may make false negatives but never false positives.

Repeating the algorithm independently reduces the error probability exponentially.

\subsection{The Class \texorpdfstring{$co$-$RP$}{co-RP}}

\paragraph{Definition ($co$-$RP$).}

A problem belongs to $co$-$RP$ if there exists a randomized polynomial-time algorithm satisfying:

\begin{itemize}
\item If the correct answer is \texttt{YES}, the algorithm always outputs \texttt{YES}.
\item If the correct answer is \texttt{NO}, the algorithm outputs \texttt{NO} with probability at least $\frac12$.
\end{itemize}

Such algorithms may make false positives but never false negatives.

\subsection{The Class \texorpdfstring{$ZPP$}{ZPP}}

\paragraph{Definition (Zero-Error Probabilistic Polynomial Time).}

A problem belongs to $ZPP$ if there exists a randomized algorithm that

\begin{itemize}
\item never outputs an incorrect answer,
\item has expected polynomial running time.
\end{itemize}

An important characterization is

\[
ZPP=RP\cap co\text{-}RP.
\]

\subsection{The Class \texorpdfstring{$BPP$}{BPP}}

\paragraph{Definition (Bounded-Error Probabilistic Polynomial Time).}

A decision problem belongs to $BPP$ if there exists a randomized polynomial-time algorithm whose probability of producing the correct answer is at least $\frac23$ on every input.

The constants $\frac23$ and $\frac13$ are arbitrary. By repeating the algorithm several times and taking a majority vote, the error probability can be reduced exponentially.

Thus $BPP$ captures problems that can be solved efficiently with an arbitrarily small probability of error.

\section{Probabilistic Polynomial-Time Algorithms}

The term \emph{Probabilistic Polynomial Time} (PPT) is frequently used in randomized algorithms and cryptography.

\paragraph{Definition (PPT Algorithm).}

A PPT algorithm is a randomized algorithm that

\begin{enumerate}
\item uses random bits during execution,
\item runs in polynomial time.
\end{enumerate}

Unlike classes such as $RP$ or $BPP$, the term PPT refers to a model of computation rather than a class of decision problems.

\paragraph{Examples.}

\begin{itemize}
\item Randomized QuickSort.
\item Miller--Rabin primality testing.
\item Karger's minimum cut algorithm.
\end{itemize}

All of these are PPT algorithms.

\section{Relationships Between the Classes}

The following inclusions are known:

\[
P\subseteq ZPP\subseteq RP\subseteq BPP.
\]

Similarly,

\[
P\subseteq ZPP\subseteq co\text{-}RP\subseteq BPP.
\]

Moreover,

\[
ZPP=RP\cap co\text{-}RP.
\]

We also know that

\[
P\subseteq NP,
\qquad
P\subseteq co\text{-}NP.
\]

Many other relationships remain unknown. In particular, we do not know whether

\[
P=NP,
\qquad
NP=co\text{-}NP,
\qquad
BPP=P.
\]

These questions lie at the heart of modern complexity theory.

\section*{Summary}

\begin{itemize}
\item $P$ contains efficiently solvable decision problems.
\item $NP$ contains problems whose \texttt{YES} instances can be verified efficiently.
\item $co$-$NP$ contains problems whose \texttt{NO} instances can be verified efficiently.
\item NP-complete problems are the hardest known problems in $NP$.
\item $RP$, $co$-$RP$, $ZPP$, and $BPP$ model different forms of randomized computation.
\item A PPT algorithm is simply a randomized polynomial-time algorithm.
\item The relationship between many of these classes remains unknown.
\end{itemize}
